\documentclass[aspectratio=1610,10pt]{beamer}

\IfFileExists{beamerthemewildcat.sty}{%
  \usetheme[
    pattern=triangle,
    color=nudarkblue,
    font=poppins,
    numbering=progressbar,
    sectionpage=on,
    agendastyle=pattern
  ]{wildcat}%
}{%
  \usetheme{Madrid}
  \usecolortheme{default}
  \setbeamertemplate{footline}[frame number]
}
\setbeamertemplate{navigation symbols}{}

\usepackage[T1]{fontenc}
\usepackage{amsmath,amssymb}
\usepackage{booktabs}
\usepackage{graphicx}
\usepackage{hyperref}
\usepackage{xcolor}

\graphicspath{{../}{./}}

\newcommand{\prr}{\mathrm{PRR}}
\newcommand{\pir}{\mathrm{PIR}}
\newcommand{\cbr}{\mathrm{CBR}}
\newcommand{\tx}{P_{\mathrm{tx}}}
\newcommand{\smallpath}[1]{\scriptsize\texttt{\detokenize{#1}}}

\title[DeePC for NR-V2X]{Data-Enabled Predictive Control for NR-V2X Sidelink Transmit-Power Adaptation}
\subtitle{Final presentation before paper drafting}
\author{Bim}
\institute{ns-3 NR-V2X / NGSIM US-101 Study}
\date{June 2026}

\IfFileExists{logo-northwestern.pdf}{%
  \titlegraphic{\includegraphics[scale=0.25]{logo-northwestern.pdf}}%
}{}

\begin{document}

\begin{frame}
  \titlepage
\end{frame}

% ============================================================
\section{Problem and Research Question}

\begin{frame}{Motivation}
  \begin{itemize}
    \item Dense V2X traffic creates time-varying sidelink contention, interference, and channel occupancy.
    \item Fixed communication settings are simple, but they do not react to measured network state.
    \item The goal is to test whether measured input-output data can support closed-loop NR-V2X parameter adaptation.
    \item This presentation focuses on a Tx-power-only DeePC controller with fixed 10 Hz beaconing.
  \end{itemize}
\end{frame}

\begin{frame}{Research Question}
  \begin{block}{Main question}
    Can Data-Enabled Predictive Control (DeePC) adapt NR-V2X sidelink transmit power online using measured PRR, PIR, CBR, and traffic-density context?
  \end{block}

  \vspace{0.5em}
  \begin{itemize}
    \item Controlled input: transmit power $\tx$.
    \item Fixed setting: beacon interval $T_b=0.1$ s, i.e., 10 Hz.
    \item Outputs: awareness-range PRR, PIR, and CBR.
    \item Context/disturbance: active vehicle count in the evaluated core region.
  \end{itemize}
\end{frame}

% ============================================================
\section{Simulation Setup}

\begin{frame}{Simulation Scenario}
  \begin{table}
    \centering
    \scriptsize
    \begin{tabular}{ll}
      \toprule
      Item & Value \\
      \midrule
      Simulator & ns-3 NR-V2X sidelink \\
      Mobility data & NGSIM US-101 processed highway trajectories \\
      Simulation duration & 300 s target campaign \\
      Current closed-loop run & partial run from 0.5 s to 162.0 s \\
      Warmup / cooldown & 10 s / 10 s \\
      KPI sample time & 0.5 s \\
      KPI window & 1.0 s \\
      Awareness range & 200 m \\
      Core filter & enabled, longitudinal core $y \in [120,415.136]$ m \\
      Logical vehicles loaded & 852 \\
      Physical UE pool & 250 UEs, dynamic active assignment \\
      Peak simultaneous active vehicles & 179 \\
      Peak active core vehicles & 118 \\
      \bottomrule
    \end{tabular}
  \end{table}
\end{frame}

\begin{frame}{Radio and CAM Parameters}
  \begin{table}
    \centering
    \scriptsize
    \begin{tabular}{ll}
      \toprule
      Parameter & Value \\
      \midrule
      Carrier frequency & 5.9 GHz \\
      Bandwidth & 10 MHz \\
      MCS & 6 \\
      Max CAM encoded size & 300 bytes \\
      CAM generation & fixed periodic, ETSI CAM trigger disabled \\
      Beacon interval & $0.1$ s fixed for closed-loop evaluation \\
      Sidelink sensing & enabled \\
      HARQ & disabled in current closed-loop campaign \\
      Subchannel size & 20 RBs \\
      Selection window & 20 slots \\
      Reservation period & 100 ms \\
      Sensing $T_1/T_2$ & 2 / 33 slots \\
      PSSCH RSRP threshold & -128 dBm \\
      Tx-power constraint & $10 \leq \tx \leq 23$ dBm \\
      \bottomrule
    \end{tabular}
  \end{table}
\end{frame}

% ============================================================
\section{Data and DeePC Formulation}

\begin{frame}{Measured KPIs}
  \begin{align*}
    \prr &= \frac{N_{\mathrm{successful}}}{N_{\mathrm{eligible}}}, \\
    y_k &= [\prr_k,\ \pir_k,\ \cbr_k]^\top .
  \end{align*}

  \begin{itemize}
    \item PRR uses eligible transmitter-receiver opportunities within 200 m at transmission time.
    \item PIR measures packet inter-reception time at receivers.
    \item CBR is the PHY busy-time fraction over the 1 s KPI window.
    \item KPIs are sampled every 0.5 s and filtered to the central corridor core.
  \end{itemize}
\end{frame}

\begin{frame}{Open-Loop Training Data}
  \begin{table}
    \centering
    \scriptsize
    \begin{tabular}{ll}
      \toprule
      Item & Value \\
      \midrule
      Training run & run 001 \\
      Original data path & \smallpath{01_training/prbs_training_250veh_run001/kpi_timeseries.csv} \\
      Clean evaluation rows & 561 \\
      Time range used & 10.0 s to 290.0 s \\
      Input for current controller & $\tx$ only \\
      Outputs & $\prr$, $\pir$, $\cbr$ \\
      Context & active vehicle count, beacon interval \\
      Counterfactual fixed-beacon CSV & \smallpath{kpi_timeseries_tb0p1_for_deepc.csv} \\
      \bottomrule
    \end{tabular}
  \end{table}

  \vspace{0.4em}
  \alert{Important caveat:} the fixed-$0.1$ training CSV overwrites only the beacon-interval column. KPI outputs still come from the original PRBS input run.
\end{frame}

\begin{frame}{Hankel Dataset}
  Sampling time is $\Delta t=0.5$ s.

  \begin{table}
    \centering
    \scriptsize
    \begin{tabular}{lll}
      \toprule
      Quantity & Samples & Time \\
      \midrule
      Past horizon $T_{\mathrm{ini}}$ & 20 & 10 s \\
      Prediction horizon $N$ & 10 & 5 s \\
      Control interval & 1 sample & 0.5 s \\
      \bottomrule
    \end{tabular}
  \end{table}

  \begin{table}
    \centering
    \scriptsize
    \begin{tabular}{ll}
      \toprule
      Hankel block & Dimension \\
      \midrule
      $U_p$ & $20 \times 307$ \\
      $U_f$ & $10 \times 307$ \\
      $Y_p$ & $60 \times 307$ \\
      $Y_f$ & $30 \times 307$ \\
      $D_p$ & $40 \times 307$ \\
      $D_f$ & $20 \times 307$ \\
      \bottomrule
    \end{tabular}
  \end{table}
\end{frame}

\begin{frame}{DeePC Signals}
  \begin{align*}
    u_k &= \tx_k, \\
    y_k &= [\prr_k,\ \pir_k,\ \cbr_k]^\top, \\
    d_k &= [n_{\mathrm{active},k},\ T_{b,k}]^\top .
  \end{align*}

  \begin{itemize}
    \item $n_{\mathrm{active}}$ is the active core vehicle count.
    \item $T_b$ is treated as fixed context, not a controlled input.
    \item Future context forecast:
      \begin{itemize}
        \item active vehicle count: hold-last forecast;
        \item beacon interval: fixed at $0.1$ s for all prediction steps.
      \end{itemize}
  \end{itemize}
\end{frame}

\begin{frame}{DeePC Optimization}
  At every 0.5 s control step, solve for the behavioral coefficient vector $g$:

  \begin{align*}
    \min_g \quad
    & \|U_f g\|_R^2 + \|Y_f g-r\|_Q^2 + \lambda_g \|g\|_2^2 \\
    \mathrm{s.t.}\quad
    & U_p g = u_{\mathrm{ini}}, \\
    & Y_p g = y_{\mathrm{ini}}, \\
    & D_p g = d_{\mathrm{ini}}, \\
    & D_f g = d_f, \\
    & 10 \leq U_f g \leq 23 .
  \end{align*}

  Only the first predicted Tx-power value is applied, then the problem is solved again at the next control step.
\end{frame}

\begin{frame}{Optimization Weights}
  \begin{table}
    \centering
    \scriptsize
    \begin{tabular}{ll}
      \toprule
      Parameter & Value \\
      \midrule
      PRR reference & 0.90 \\
      PIR reference & 0.10 s \\
      CBR reference & 0.35 \\
      Input penalty $R$ & 0.002 repeated over 10 steps \\
      Regularization $\lambda_g$ & 10.0 \\
      PRR-only weights & $q_{\prr}=40,\ q_{\pir}=0,\ q_{\cbr}=0$ \\
      CBR-only weights & $q_{\prr}=0,\ q_{\pir}=0,\ q_{\cbr}=8$ \\
      PIR-only weights & $q_{\prr}=0,\ q_{\pir}=3,\ q_{\cbr}=0$ \\
      Solver & CVXPY with OSQP \\
      Residual tolerance & $10^{-5}$ \\
      \bottomrule
    \end{tabular}
  \end{table}
\end{frame}

% ============================================================
\section{Closed-Loop Implementation}

\begin{frame}{Closed-Loop Architecture}
  \begin{enumerate}
    \item ns-3 collects the latest 20 samples of $u$, $y$, and $d$.
    \item ns-3 writes a JSON request to the file bridge.
    \item Python/CVXPY solves the DeePC QP using OSQP.
    \item Python writes predicted 10-step trajectories and the first control action.
    \item ns-3 applies only the first Tx-power value for the next 0.5 s.
  \end{enumerate}

  \vspace{0.5em}
  \small
  Current run path:
  \smallpath{data/output/final/ieee_250veh_deepc_campaign_01/02_runs/deepc_count_tb0p1_prr_only/run_931}
\end{frame}

\begin{frame}{Closed-Loop Command}
  \scriptsize
  \begin{block}{PRR-only fixed-10 Hz DeePC run}
  \texttt{DATASET\_DIR=data/output/final/ieee\_250veh\_deepc\_campaign\_01/01\_training/deepc\_dataset\_txpower\_count\_tb0p1\_dt0p5\_n10 \textbackslash}\\
  \texttt{METHOD=deepc\_count\_tb0p1\_prr\_only \textbackslash}\\
  \texttt{DEEPC\_CONTROLLER\_ARGS="--q-prr 40 --q-pir 0 --q-cbr 0" \textbackslash}\\
  \texttt{data/output/final/ieee\_250veh\_deepc\_campaign\_01/commands/run\_deepc\_python\_txpower.sh 931}
  \end{block}

  \vspace{0.5em}
  \begin{itemize}
    \item Target duration: 300 s.
    \item Current completed partial result: 162 s.
    \item Plot folder: \smallpath{run_931/plots}
  \end{itemize}
\end{frame}

% ============================================================
\section{Results}

\begin{frame}{Main Result Summary}
  \begin{table}
    \centering
    \scriptsize
    \begin{tabular}{lrrrrrr}
      \toprule
      Method & Time (s) & PRR & PIR (s) & CBR & Mean $\tx$ & Mean $T_b$ \\
      \midrule
      Fixed 10 Hz & 161.0 & 0.541 & 0.151 & 0.386 & 20.00 & 0.10 \\
      DeePC PRR-only & 162.0 & 0.540 & 0.151 & 0.386 & 18.10 & 0.10 \\
      Fixed 5 Hz & 119.5 & 0.467 & 0.256 & 0.273 & 20.00 & 0.20 \\
      \bottomrule
    \end{tabular}
  \end{table}

  \vspace{0.5em}
  \begin{itemize}
    \item DeePC achieved fixed-10 Hz-like PRR/PIR/CBR over the completed interval.
    \item Mean applied Tx power was reduced from 20 dBm to about 18.10 dBm.
    \item This is a feasibility result, not yet a definitive reliability improvement.
  \end{itemize}
\end{frame}

\begin{frame}{Controller Health}
  \begin{table}
    \centering
    \scriptsize
    \begin{tabular}{lr}
      \toprule
      Quantity & Value \\
      \midrule
      KPI rows & 324 \\
      Control updates & 285 \\
      Controller success rate & 100\% \\
      Mean solve time & 0.0199 s \\
      Max solve / bridge wait time & 0.167 s / 0.202 s \\
      Mean applied Tx power & 18.10 dBm \\
      Mean beacon interval & 0.1 s \\
      Max equality residuals & 0.0 in logged run \\
      Prediction rows logged & 2850 \\
      \bottomrule
    \end{tabular}
  \end{table}
\end{frame}

\begin{frame}{Closed-Loop KPI Time Series}
  \centering
  \includegraphics[width=0.92\linewidth]{data/output/final/ieee_250veh_deepc_campaign_01/02_runs/deepc_count_tb0p1_prr_only/run_931/plots/kpi_timeseries.png}
\end{frame}

\begin{frame}{Applied Control}
  \centering
  \includegraphics[width=0.92\linewidth]{data/output/final/ieee_250veh_deepc_campaign_01/02_runs/deepc_count_tb0p1_prr_only/run_931/plots/controls.png}
\end{frame}

\begin{frame}{Prediction Trajectories}
  \centering
  \includegraphics[width=0.92\linewidth]{data/output/final/ieee_250veh_deepc_campaign_01/02_runs/deepc_count_tb0p1_prr_only/run_931/plots/prediction_trajectories_10_step.png}
\end{frame}

\begin{frame}{Prediction Error and Timing}
  \begin{columns}
    \begin{column}{0.49\linewidth}
      \centering
      \includegraphics[width=\linewidth]{data/output/final/ieee_250veh_deepc_campaign_01/02_runs/deepc_count_tb0p1_prr_only/run_931/plots/prediction_error_by_horizon.png}
    \end{column}
    \begin{column}{0.49\linewidth}
      \centering
      \includegraphics[width=\linewidth]{data/output/final/ieee_250veh_deepc_campaign_01/02_runs/deepc_count_tb0p1_prr_only/run_931/plots/controller_timing.png}
    \end{column}
  \end{columns}
\end{frame}

% ============================================================
\section{Discussion and Next Steps}

\begin{frame}{Discussion}
  \begin{itemize}
    \item The closed-loop bridge is working: OSQP returns feasible controls at every logged control step.
    \item The controller maintains fixed-10 Hz-like PRR/PIR/CBR while using lower mean Tx power.
    \item The current result is conservative: PRR does not improve clearly over fixed 10 Hz, but power is reduced.
    \item Context helps make the DeePC problem density-conditioned, but hold-last vehicle-count forecasting is simple.
    \item The fixed-$0.1$ Hankel dataset is counterfactual because it reuses outputs from a PRBS beacon-interval training run.
  \end{itemize}
\end{frame}

\begin{frame}{What Is Scientifically Strong Already?}
  \begin{itemize}
    \item Realistic microscopic mobility from NGSIM US-101.
    \item Dynamic physical-UE pool avoids creating one NR UE per logical vehicle.
    \item Closed-loop ns-3-to-Python controller runs online, not post-processed offline only.
    \item DeePC optimization uses measured histories and Hankel behavior constraints.
    \item Controller diagnostics, predicted trajectories, residuals, and timing are logged.
  \end{itemize}
\end{frame}

\begin{frame}{Limitations Before Paper Submission}
  \begin{itemize}
    \item Need full 300 s completed runs for PRR-only, CBR-only, and PIR-only objectives.
    \item Need repeated seeds or repeated runs for credible comparison.
    \item Need a clean fixed-10 Hz Tx-power-excitation training run, instead of counterfactual beacon-column overwrite.
    \item Need compare against fixed 10 Hz and fixed 5 Hz over matched time windows.
    \item Need final choice on HARQ setting and align all baselines with that choice.
  \end{itemize}
\end{frame}

\begin{frame}{Draft Paper Story}
  \begin{enumerate}
    \item Build a realistic NR-V2X NGSIM simulation with packet-level PRR/PIR/CBR logging.
    \item Collect open-loop input-output training data.
    \item Build Hankel matrices for Tx-power-only DeePC.
    \item Add active vehicle count and fixed beacon interval as measured context.
    \item Run closed-loop DeePC through a file bridge with OSQP.
    \item Show feasibility, controller timing, prediction quality, and comparison against fixed-rate baselines.
  \end{enumerate}
\end{frame}

\begin{frame}{Takeaway}
  \begin{block}{Main message}
    DeePC can be integrated with an ns-3 NR-V2X sidelink simulation and can control transmit power online under realistic NGSIM traffic while maintaining fixed-10 Hz-like KPIs with lower mean transmit power.
  \end{block}

  \vspace{0.8em}
  \begin{block}{Next paper step}
    Replace the counterfactual fixed-beacon Hankel dataset with a clean fixed-10 Hz open-loop Tx-power-excitation dataset and repeat matched closed-loop baseline comparisons.
  \end{block}
\end{frame}

\end{document}
